
\documentclass{article}
\usepackage{colortbl}
\usepackage{makecell}
\usepackage{multirow}
\usepackage{supertabular}

\begin{document}

\newcounter{utterance}

\centering \large Interaction Transcript for game `cladder', experiment `full\_v1.5\_default', episode 3258 with selfrefine{-}llama{-}3.1{-}8b{-}instruct.
\vspace{24pt}

{ \footnotesize  \setcounter{utterance}{1}
\setlength{\tabcolsep}{0pt}
\begin{supertabular}{c@{$\;$}|p{.15\linewidth}@{}p{.15\linewidth}p{.15\linewidth}p{.15\linewidth}p{.15\linewidth}p{.15\linewidth}}
    \# & \multicolumn{2}{c}{Player} && \multicolumn{2}{c}{Game Master} \\
    \hline

    \theutterance \stepcounter{utterance}  
    & & & \multicolumn{4}{p{0.6\linewidth}}{
        \cellcolor[rgb]{0.9,0.9,0.9}{
            \makecell[{{p{\linewidth}}}]{
                \texttt{\tiny{[P1$\langle$GM]}}
                \texttt{You are an expert in causal inference. The following question is not a typical commonsense query, but rather a meticulously designed question created by a professor specializing in causal inference, intended to assess the students' mastery of the course content.} \\
\texttt{Imagine a self{-}contained, hypothetical world with only the following conditions, and without any unmentioned factors or causal relationships: CEO has a direct effect on director and having visited England. Having visited England has a direct effect on employee. Director has a direct effect on employee. The overall probability of having visited England is 19\%. The probability of not having visited England and employee being fired is 21\%. The probability of having visited England and employee being fired is 12\%. Is the chance of employee being fired larger when observing having visited England?} \\
\texttt{Start your answer with 'Yes' or 'No', followed by additional reasoning or evidence to support your explanation.} \\
            }
        }
    }
     \\ \\

    \theutterance \stepcounter{utterance}  
    & & \multicolumn{4}{p{0.6\linewidth}}{
        \cellcolor[rgb]{0.95,0.95,0.95}{
            \makecell[{{p{\linewidth}}}]{
                \texttt{\tiny{[P1$\leftrightarrow$P]}}
                \texttt{{[}DRAFT{]}} \\
\texttt{No.} \\
\\ 
\texttt{To determine whether the chance of employee being fired is larger when observing having visited England, we need to calculate the probability of employee being fired given having visited England (P(Fired|Visited England)) and compare it to the probability of employee being fired in the absence of having visited England (P(Fired|Not Visited England)).} \\
\\ 
\texttt{We are given the following probabilities:} \\
\\ 
\texttt{{-} P(Visited England) = 19\%} \\
\texttt{{-} P(Not Visited England) = 81\%} \\
\texttt{{-} P(Fired|Not Visited England) = 21\%} \\
\texttt{{-} P(Fired|Visited England) = 12\%} \\
\\ 
\texttt{We can calculate the probability of employee being fired given having visited England using the law of total probability:} \\
\\ 
\texttt{P(Fired|Visited England) = P(Fired|Visited England) * P(Visited England) / P(Visited England)} \\
\texttt{However, we can also calculate P(Fired|Not Visited England) * P(Not Visited England) / P(Not Visited England) to get the same result.} \\
\\ 
\texttt{To compare the two probabilities, we can calculate the odds ratio of employee being fired given having visited England versus not having visited England:} \\
\\ 
\texttt{Odds Ratio = P(Fired|Visited England) / P(Fired|Not Visited England) * P(Not Visited England) / P(Visited England)} \\
\\ 
\texttt{Plugging in the values, we get:} \\
\\ 
\texttt{Odds Ratio = 12\% / 21\% * 81\% / 19\% ≈ 0.57} \\
\\ 
\texttt{Since the odds ratio is less than 1, it indicates that the chance of employee being fired is smaller when observing having visited England.} \\
\\ 
\texttt{Therefore, the answer is No.} \\
            }
        }
    }
    & \\ \\

    \theutterance \stepcounter{utterance}  
    & & \multicolumn{4}{p{0.6\linewidth}}{
        \cellcolor[rgb]{0.95,0.95,0.95}{
            \makecell[{{p{\linewidth}}}]{
                \texttt{\tiny{[P1$\leftrightarrow$P]}}
                \texttt{{[}SELF{-}FEEDBACK{]}} \\
\texttt{FORMAT: ok} \\
\texttt{RULES: ok} \\
\texttt{STRATEGY: problem} \\
\texttt{REASON: The answer incorrectly concludes that the chance of employee being fired is smaller when observing having visited England, despite the odds ratio indicating the opposite.} \\
            }
        }
    }
    & \\ \\

    \theutterance \stepcounter{utterance}  
    & \multicolumn{4}{p{0.6\linewidth}}{
        \cellcolor[rgb]{0.561,0.737,0.561}{
            \makecell[{{p{\linewidth}}}]{
                \texttt{\tiny{[P1$\rangle$GM]}}
                \texttt{No.} \\
\\ 
\texttt{To determine whether the chance of employee being fired is larger when observing having visited England, we need to calculate the probability of employee being fired given having visited England (P(Fired|Visited England)) and compare it to the probability of employee being fired in the absence of having visited England (P(Fired|Not Visited England)).} \\
\\ 
\texttt{We are given the following probabilities:} \\
\\ 
\texttt{{-} P(Visited England) = 19\%} \\
\texttt{{-} P(Not Visited England) = 81\%} \\
\texttt{{-} P(Fired|Not Visited England) = 21\%} \\
\texttt{{-} P(Fired|Visited England) = 12\%} \\
\\ 
\texttt{We can calculate the probability of employee being fired given having visited England using the law of total probability:} \\
\\ 
\texttt{P(Fired|Visited England) = P(Fired|Visited England) * P(Visited England) / P(Visited England)} \\
\texttt{However, we can also calculate P(Fired|Not Visited England) * P(Not Visited England) / P(Not Visited England) to get the same result.} \\
\\ 
\texttt{To compare the two probabilities, we can calculate the odds ratio of employee being fired given having visited England versus not having visited England:} \\
\\ 
\texttt{Odds Ratio = P(Fired|Visited England) / P(Fired|Not Visited England) * P(Not Visited England) / P(Visited England)} \\
\\ 
\texttt{Plugging in the values, we get:} \\
\\ 
\texttt{Odds Ratio = 12\% / 21\% * 81\% / 19\% ≈ 0.57} \\
\\ 
\texttt{Since the odds ratio is less than 1, it indicates that the chance of employee being fired is actually larger when observing having visited England, but the odds ratio is not the correct measure to determine this. The correct measure is the probability of employee being fired given having visited England (P(Fired|Visited England)) which is 12\%, and the probability of employee being fired in the absence of having visited England (P(Fired|Not Visited England)) which is 21\%. The probability of employee being fired given having visited England is smaller than the probability of employee being fired in the absence of having visited England.} \\
\\ 
\texttt{Therefore, the answer is No.} \\
            }
        }
    }
    & & \\ \\

    \theutterance \stepcounter{utterance}  
    & & & \multicolumn{2}{p{0.3\linewidth}}{
        \cellcolor[rgb]{0.9,0.9,0.9}{
            \makecell[{{p{\linewidth}}}]{
                \texttt{\tiny{[GM$|$GM]}}
                \texttt{No} \\
            }
        }
    }
    & & \\ \\

    \theutterance \stepcounter{utterance}  
    & & & \multicolumn{2}{p{0.3\linewidth}}{
        \cellcolor[rgb]{0.9,0.9,0.9}{
            \makecell[{{p{\linewidth}}}]{
                \texttt{\tiny{[GM$|$GM]}}
                \texttt{yes} \\
            }
        }
    }
    & & \\ \\

    \theutterance \stepcounter{utterance}  
    & & & \multicolumn{2}{p{0.3\linewidth}}{
        \cellcolor[rgb]{0.9,0.9,0.9}{
            \makecell[{{p{\linewidth}}}]{
                \texttt{\tiny{[GM$|$GM]}}
                \texttt{game\_result = LOSE} \\
            }
        }
    }
    & & \\ \\

\end{supertabular}
}

\end{document}
